\section{Boolean message passing}\label{sec:booleanMessagePassing}

We investigate message passing schemes with respect to the Boolean semiring $\booleansemiring$.
To this end, we derive the schemes based on properties of the Boolean semiring and then show their application to sparsify generic contractions.

\subsection{Properties of the Boolean semiring}

The main property of contractions in the Boolean semiring is the invariance under adding subcontractions, which we state in the next theorem.

\begin{theorem}\label{the:booleanContractionInvariance}
    Let $\extnet$ be an tensor network on an arbitrary hypergraph $\graph=(\nodes,\edges)$ with respect to the Boolean semiring.
    For any subsets $\secnodes\subset\nodes$ and $\secedges\subset\edges$ we have
    \begin{align*}
        \srcontractionofwrt{\extnet\cup
            \{\srcontractionofwrt{\{\hypercoreofat{\edge}{\catvariableof{\edge}}\wcols\edge\in\secedges\}}{\catvariableof{\secnodes}}{\booleansemiring}\}
        }{\catvariableof{\nodes}}{\booleansemiring}
        = \srcontractionofwrt{\extnet}{\catvariableof{\nodes}}{\booleansemiring} \, .
    \end{align*}
\end{theorem}
\begin{proof}
    Coordinatewise:
    Only invariant, if subcontraction evaluates to $0$.
    But then by monotonocity of subcontractions also the whole contraction evaluates to $0$, therefore invariant.
\end{proof}

\subsection{Application on sum-product contractions}

\theref{the:booleanContractionInvariance} can be generalized to contractions on more general semirings sharing the elements $0,1$ with $\booleansemiring$.
In such way we can switch to message passing schemes in the Boolean semiring to sparsify more generic tensor networks.

\subsection{Generic support tensors}

%% Generic support tensor
Even more generic, we define the support operator on a generic tensor with respect to a semiring $\gensemiring$ as
\begin{align*}
    \srnonzerocirc\hypercoreat{\indexedshortcatvariables}
    = \begin{cases}
          \srzero & \ifspace \hypercoreat{\indexedshortcatvariables} = \srzero \\
          \srone & \elsetext \, .
    \end{cases}
\end{align*}
We notice that the support operator leaves any tensor over semirings with $S=\{\srzero,\srone\}$ invariant.
This is the case for the Boolean semiring $\booleansemiring$ and the Galois field $\galoissemiring$.

\begin{lemma}\label{lem:supportInvarianceGeneric}
    Any tensor is invariant under contraction with its support tensor, that is
    \begin{align*}
        \hypercorewith
        = \srcontractionofwrt{
            \hypercorewith,\srnonzerocirc\hypercorewith
        }{\shortcatvariables}{\gensemiring} \, .
    \end{align*}
\end{lemma}
\begin{proof}
    For any coordinate $s\neq\srzero$ the support operator returns $\srone$, and by the multiplicative invariance property we have $s \srmul \srone = s$.
    If $s=\srzero$ the support operator returns $\srzero$ and we have $\srzero\srmul\srzero=\srzero$.
\end{proof}

\begin{definition}
    A semiring $\gensemiring$ is called \emph{inverse-free}, if for any $s_0,s_1\in S$ with $s_0\srplus s_1=\srzero$ we have $s_0=s_1=\srzero$.
\end{definition}

% Examples
Note that $([0,\infty),+,\cdot,0,1)$ is inverse-free, while $(\rr,+,\cdot,0,1)$ is not.
The Boolean semiring $\booleansemiring$ is inverse-free, but the Galois field $\galoissemiring$ is not.

\begin{lemma}\label{lem:supportDecompositionGeneric}
    The support of a contraction obeys
    \begin{align*}
        \srnonzerocirc\srcontractionofwrt{\extnet}{\catvariableof{\secnodes}}{\gensemiring}
        \prec \srcontractionofwrt{
            \{\srnonzerocirc\hypercoreofat{\edge}{\catvariableof{\edge} \wcols \edgein}\}
        }{\catvariableof{\secnodes}}{\booleansemiring} \, .
    \end{align*}
    Here the Boolean semiring is defined with the neutral elements $\srzero,\srone$ of the semiring $\gensemiring$, and the inequality denotes partial ordering of Boolean tensors.
    If the semiring is inverse free, then the above inequality holds straight.
\end{lemma}
\begin{proof}
    To show the first claim, let $\catindexof{\secnodes}$ be an index such that
    \begin{align*}
        \srcontractionofwrt{
            \{\srnonzerocirc\hypercoreofat{\edge}{\catvariableof{\edge}} \wcols \edgein\}
        }{\indexedcatvariableof{\secnodes}}{\booleansemiring}
        = \srzero \, .
    \end{align*}
    It suffices to show that for this index we have $\srcontractionofwrt{\extnet}{\indexedcatvariableof{\secnodes}}{\gensemiring}=\srzero$.
    Since the Boolean semiring is inverse-free we have for any extension $\catindexof{\nodes}$ of $\catindexof{\secnodes}$ to $\nodes$ that
    \begin{align*}
        \srcontractionofwrt{
            \{\srnonzerocirc\hypercoreofat{\edge}{\catvariableof{\edge} \wcols \edgein}\}
        }{\indexedcatvariableof{\nodes}}{\booleansemiring}
        = \srzero \, .
    \end{align*}
    We therefore find $\edgein$ such that for the restriction $\catindexof{\edge}$ of $\catindexof{\nodes}$ to $\edge$ we have
    \begin{align*}
        \hypercoreofat{\edge}{\indexedcatvariableof{\edge}}
        = \srzero
    \end{align*}
    and thus
    \begin{align*}
        \srcontractionofwrt{\extnet}{\indexedcatvariableof{\nodes}}{\gensemiring}
        = \srzero \, .
    \end{align*}
    It follows that also
    \begin{align*}
        \srcontractionofwrt{\extnet}{\indexedcatvariableof{\secnodes}}{\gensemiring}
        = \srzero \, .
    \end{align*}

    The second claim, namely the vise versa inequality, can be shown analogously using that $\gensemiring$ is inverse-free.
\end{proof}

% Non-equality example
As an example, where the inequality in \lemref{lem:supportDecompositionGeneric} does not hold straight, consider the contraction of
\begin{align*}
    \hypercoreofat{0}{\catvariableof{0}}
    \coloredmatrixof{
        1 \\ -1
    }{\catvariableof{0}}
    \quad \text{and} \quad
    \hypercoreofat{1}{\catvariableof{0},\catvariableof{1}}
    \coloredmatrixof{
        1 & 1 \\ 1 & 1
    }{\catvariableof{0},\catvariableof{1}} \, .
\end{align*}
The support is the trivial tensor in both cases, but $\contractionof{\hypercoreofat{0}{\catvariableof{0}},\hypercoreofat{1}{\catvariableof{0},\catvariableof{1}}}{\catvariableof{1}}=\onesat{\catvariableof{1}}$.
When we restrict to the non-negative real numbers in the sum-product semiring, the semiring is inverse free and \lemref{lem:supportDecompositionGeneric} would guarantee equality.

Boolean message passing can therefore be applied to sparsify generic contractions to arbitrary semirings:
\begin{itemize}
    \item Add the support of the contraction (\lemref{lem:supportInvarianceGeneric})
    \item Decompose the support into local network (\lemref{lem:supportDecompositionGeneric})
    \item Perform Boolean message passing to simplify the support
\end{itemize}
